\section{The algorithm}\label{sec:method}

Given points $p_0,\dots,p_{n-1}$ in the plane with distance matrix $D_{ij}=\lVert p_i-p_j\rVert_2$
(rounded to the nearest integer for TSPLIB \code{EUC\_2D} instances, as the library prescribes), BGMA
returns a Hamiltonian cycle. A tour is an array of city indices; for de-duplication it is put in a
canonical form that starts at city~0 and takes the lexicographically smaller of the two directions.
Preprocessing sorts, for every city, all other cities by distance ($O(n^2\log n)$) and keeps a
truncated candidate list of the $m=\max(10,\operatorname{round}(n/10))$ nearest.

BGMA runs one or more independent \emph{units}. A unit (Algorithm~\ref{alg:unit}) has five parts.

\paragraph{Beam islands.}
A beam state is a partial path with a start city, a current city, a visited set and a length. An island
starts from two \emph{poles}: the farthest pair of cities for the first island, and two seeded random
cities for the others. Each pole seeds its own start state. At each of the $n-1$ extension steps
every state proposes at most four successors: the nearest unvisited neighbour of the current city,
the second-nearest unvisited neighbour of the previous city, the third-nearest unvisited neighbour of
the city before that, and the second-nearest unvisited neighbour of the current city. Successors
are merged on a Zobrist hash of (visited set, last city, start city), keeping the shorter path.
They are scored by $f=\text{length}+D_{\text{last},\text{start}}$. The $K$ best survive. The $K$
completed tours (the \emph{finals}) are all passed on. The beam is a cheap generator of diverse,
structurally reasonable tours, in the spirit of filtered beam search~\citep{ow1988}; it is not where
the tour quality comes from.

\paragraph{Fast screen.}
Every beam final and every child is improved by a candidate-list local search: 2-opt~\citep{croes1958}
over the $m$-nearest lists with don't-look bits and a work queue~\citep{bentley1992}, then Or-opt
relocation of segments of length $\le 3$~\citep{or1976} (both orientations, insertion next to a
candidate neighbour), then 2-opt again. Iteration caps bound the work: at most $200n$ queue pops for 2-opt, and at most 20 accepted Or-opt
moves, each followed by a rescan from the start of the tour. The cost is close to linear in $n$ in
practice, but the result is not in general a local optimum: on large instances the Or-opt cap binds,
and the full polish below does the remaining work.

\paragraph{Elite pool and recombination.}
Screened tours are de-duplicated by canonical form, sorted by length and truncated to the best $e$.
A generation creates $c$ children. The first parent is drawn uniformly from the best $\min(4,|P|)$
tours, the second uniformly from the whole pool. The child is built by edge recombination
(ERX)~\citep{whitley1989}: from a random start city, repeatedly move to the unvisited neighbour in the
union of the parents' edges that has the fewest unvisited neighbours itself (ties: the shorter
edge), and jump to the nearest unvisited city when no parental edge is available.

\paragraph{Perturbation.}
Two independent events may then modify the child. With probability \code{kickRate} it receives one
double-bridge move~\citep{martin1991}. Independently, with probability \code{annealRate}, it receives a
short simulated anneal~\citep{kirkpatrick1983}. The anneal uses candidate-list 2-opt proposals and
Metropolis acceptance. The temperature starts at $T_0=\max(1,0.02n)$ and is multiplied by $0.99998$
per proposal. After $50n$ proposals without a new best it applies a double bridge and reheats to
$\max(T,0.4T_0)$. The anneal stops after $\min(A,\max(1000,200n))$ proposals, where $A$ is the level's
cap, and returns the best tour it saw. $T_0$ is in absolute length units and is not normalised by the
instance scale; we return to this in Section~\ref{sec:limits}.

\paragraph{Selective full polish.}
After screening, the children are sorted and only the best $\min(q,c)$ receive the expensive local search:
first-improvement 2-opt over all $O(n^2)$ pairs followed by Or-opt with segments up to $L_p$ cities.
The scan restarts after every accepted move and stops at a true local optimum for that
neighbourhood. The same polish is applied to the best $\min(q,|P|)$ tours of the initial pool and, if
\code{finalPolish}~$>0$, to that many pool tours at the end. Restricting the quadratic search to
contenders is the main cost-control device: the fast screen is the filter, the full polish is the
refinement.

\paragraph{Stagnation injection.}
If the best tour has not improved for two consecutive generations, one fresh beam island with random
poles is run, its finals are screened and merged into the pool, and the counter resets.

\begin{algorithm}[t]
\caption{One BGMA unit. All random choices come from a generator seeded by (instance, seed).}
\label{alg:unit}
\begin{algorithmic}[1]
\State $S\gets$ finals of $I$ beam islands (width $K$), each improved by the fast screen
\State $P\gets$ best $e$ distinct tours of $S$;\quad polish the best $q$ tours of $P$ and merge them back
\For{$g=1,\dots,G$}
  \For{$j=1,\dots,c$}
    \State $a\gets$ uniform from the best $\min(4,|P|)$ of $P$;\quad $b\gets$ uniform from $P$
    \State $x\gets \mathrm{ERX}(a,b)$
    \State with probability \code{kickRate}: $x\gets$ double-bridge$(x)$
    \State with probability \code{annealRate} (independently): $x\gets$ anneal$(x)$
    \State $C_j\gets$ fast-screen$(x)$
  \EndFor
  \State polish the best $\min(q,c)$ of $C$;\quad $P\gets$ best $e$ distinct tours of $P\cup C$
  \If{no improvement for 2 generations} merge one screened random-pole beam island into $P$ \EndIf
\EndFor
\State polish the best \code{finalPolish} tours of $P$ and merge;\quad \Return best tour in $P$
\end{algorithmic}
\end{algorithm}

\paragraph{Restarts and the effort dial.}
A level with $R$ restarts runs $R$ units with seed offsets and keeps the shortest tour. The five levels
\lvl{low}, \lvl{mid}, \lvl{high}, \lvl{xhigh} and \lvl{max} (Table~\ref{tab:ladder}) differ only in
budgets. The top level is a portfolio: \lvl{max} runs the \lvl{xhigh} configuration once and the
\lvl{max} configuration with $R=2$, i.e.\ three independent units, and returns the shortest of the three
tours. By construction \lvl{max} is never worse than \lvl{xhigh} on the same instance and seed, and it
costs the sum of its units. No other monotonicity between levels is guaranteed, and we observe
violations (Section~\ref{sec:synthetic}). Because the units share nothing, they can run on separate
threads; the repository's worker pool does this and returns bit-identical results.

\begin{table}[t]
\centering\footnotesize
\caption{The effort ladder (\code{src/hybrid/effort\_profiles.js}). $r(x)=\max(1,\operatorname{round}(x))$.
$A$ and $p_A$ are the anneal cap and probability, $p_k$ the kick probability, $L_p$ the polish
Or-opt segment cap. \lvl{max} additionally runs one \lvl{xhigh} unit (portfolio). At most $c$
children are polished per generation, so the cap $q$ binds per generation only while $q<c$: from
$n\approx 200$, 175 and 210 on, \lvl{high}, \lvl{xhigh} and \lvl{max} polish all $c$ children.}
\label{tab:ladder}
\setlength{\tabcolsep}{2.5pt}
\begin{tabular}{@{}lccccccccccc@{}}
\toprule
level & $I$ & $K$ & $e$ & $G$ & $c$ & $q$ & final & $R$ & $A$ / $p_A$ & $p_k$ & $L_p$\\
\midrule
\lvl{low}   & 2 & $\max(12,r(n/4))$  & $\max(8,r(n/20))$  & $\max(2,r(n/50))$  & 4  & 2                 & 0 & 1 & 8k / .25  & .30 & 5\\
\lvl{mid}   & 4 & $\max(16,r(n/3))$  & $\max(16,r(n/10))$ & $\max(4,r(n/25))$  & 10 & 6                 & 2 & 1 & 20k / .35 & .35 & 5\\
\lvl{high}  & 5 & $\max(20,r(n/2.5))$& $\max(24,r(n/6))$  & $\max(6,r(n/15))$  & 14 & $\max(6,r(n/15))$ & 3 & 1 & 30k / .40 & .40 & 6\\
\lvl{xhigh} & 6 & $\max(24,r(n/2))$  & $\max(32,r(n/3))$  & $\max(8,r(n/8))$   & 18 & $\max(10,r(n/10))$& 6 & 1 & 50k / .45 & .45 & 8\\
\lvl{max}   & 8 & $\max(32,r(n/2))$  & $\max(48,r(n/2))$  & $\max(10,r(n/8))$  & 18 & $\max(10,r(n/12))$& 8 & 2 & 60k / .50 & .50 & 8\\
\bottomrule
\end{tabular}
\end{table}

\paragraph{Cost.}
Per unit the work is roughly $I$ beam islands (about $4nK$ node creations, each copying an $n$-bit visited set),
$IK+Gc$ fast screens, $\min(q,e)+G\min(q,c)+\text{final}$ full polishes, and up to $Gc\,p_A$ anneals
of at most $A$ proposals. The full polish dominates at large $n$. One scan is $O(n^2)$ and is restarted
after every accepted move, and in the upper levels $G$ grows linearly with $n$, while the polishes per
generation are capped at $c$. The budgets are
polynomial in $n$ by construction, so BGMA is a polynomial-time heuristic and carries no optimality or
approximation guarantee. The Euclidean TSP is NP-hard~\citep{papadimitriou1977}, and the result of
any such heuristic says nothing about P versus NP.

\paragraph{The parameters were tuned in-sample.}
The ladder was set by one-factor sweeps on the same generator of uniform and clustered synthetic
instances that Section~\ref{sec:synthetic} reports on. The TSPLIB instances in
Section~\ref{sec:tsplib} were not used for tuning.
